\section{Related Work}

The pipeline builds on LLM-as-judge research, but adapts it to source-grounded
legal reasoning rather than general helpfulness or preference evaluation.
MT-Bench and Chatbot Arena showed that strong LLM judges can approximate human
preferences at scale, while also surfacing biases such as position, verbosity,
and self-enhancement effects \cite{zheng2023judging}. G-Eval introduced a
rubric-like evaluation procedure using chain-of-thought-style criteria and
reported stronger human alignment for natural-language generation evaluation
\cite{liu2023geval}. Prometheus and Prometheus 2 further support the idea that
rubric-conditioned evaluator models can provide fine-grained judgments when
given reference material and score criteria \cite{kim2023prometheus,
kim2024prometheus2}.

The judge literature also makes clear why internal validation must test more
than whether an LLM can emit a score. Position-bias studies show that pairwise
or listwise judges may change preferences when candidate order changes
\cite{wang2024positionbias}. Self-preference work shows that judge models may
favor outputs from their own model family \cite{panickssery2024selfpreference}.
For legal evaluation specifically, recent benchmarks and legal judge proposals
emphasize that legal reasoning requires decomposed factual and doctrinal units,
not just general answer preference \cite{lexam2025,lemaj2025}. These findings
motivate our row-level rubric design, centroid-level judging, model-output
provenance, and explicit separation between internal validation and future
expert-reviewed legal validity.

Our pipeline differs in three ways. First, the benchmark is generated from a
legal source case rather than from a static task prompt. Second, responses are
clustered by legal reasoning path before judging, allowing many model samples to
be compressed into interpretable centroids. Third, the judge scores rubric rows
on centroids and then projects those scores to every response in the cluster.
This makes model ranking cheaper and more interpretable while preserving
row-level evidence for review.
